\section{Energy of the last untrimmed family image}

\begin{lemma}[A terminal compressed soft-energy estimate]
    \label{thm:soft_terminal_compressed_energy}
    \lean{TNLean.MPS.AreaLaw.ALVV.exists_native_soft_terminal_compressed_energy_estimate}
    \leanok
    Let $H=\sum_i h_i$ be a chain of positive local dimension, with
    $0\le h_i\le I$ and the designated nearest-neighbour supports.
    Let $G$ be its full ground space. Fix $l\le u\le v\le r\le L$,
    with $u-l,r-v\le2\ell$, and let $P,M$ be the original internal
    interaction sums on $J=[l,r)$ and $[u,v)$, respectively.
    Their actual attained minima $e_P,e_M$ are chosen before the
    center and every terminal candidate.

    Suppose $e_M-10\le p\le e_M$, $t,T_0\ge48$, and $m\ge2$.
    Define
    \begin{align}
        R      & =P-M+pI+h_{m,t}(M-pI),                                     &
        A_*    & =t(1+\log m)+4\ell+14,\label{eq:terminal_soft_operator}      \\
        \rho   & =32e^{-(T_0-48)/8}+\frac tm\left(\frac{T_0}t\right)^m,     &
        \delta & =\frac{\Delta}{(91/100)^2}.\label{eq:terminal_soft_losses}
    \end{align}
    Assume $\tau,\Delta\ge0$, $\rho\le1/16$,
    $\tau\le1/[1024(A_*+1)]$ and $\delta\le1/[100(A_*+1)]$.
    Restrict $R$ to a fixed complementary configuration to obtain
    $R_{\mathrm{nat}}$ on $\mathcal H_J$. Then
    $R=R_{\mathrm{nat}}\otimes I$ in the explicit regional tensor
    coordinates. There is an actual attained minimum $e_{\mathrm s}$,
    chosen before the target and candidate, such that
    $R_{\mathrm{nat}}-e_{\mathrm s}I\ge0$,
    $\|R-e_{\mathrm s}I\|\le A_*$ and
    $e_{\mathrm s}\le e_P\le e_{\mathrm s}+\rho$.

    Let $T$ be a full-chain space with $\dim T=\dim G$ and
    $\|(I-P_G)t_0\|^2\le\tau\|t_0\|^2$ for $t_0\in T$.
    Let $I_e:\mathbb C^d\to\mathcal H_J$ be the coordinate isometry,
    let $(A_i)_i$ be a finite operator family on $\mathcal H_J$,
    and let $Q$ be any $d\times d$ matrix. No projection assumption
    on $Q$ is imposed. Form the same untrimmed family image
    $Y=\sum_i A_i I_e\operatorname{range}Q$.
    Suppose $Y$ is $\delta$-viable for $T$ in the explicit regional
    coordinates, retaining the whole exterior of $J$.
    Then $Y$ has positive dimension, and the actual orthogonal
    compression of $R_{\mathrm{nat}}$ to $Y$ has an orthonormal
    eigenbasis with decreasing eigenvalues. Its least eigenvalue
    $\lambda$ satisfies
    \begin{align}
        e_{\mathrm s}\le\lambda\le e_{\mathrm s}+3,\qquad
        e_P-3-\rho\le\lambda-3\le e_P,\qquad
        e_P-49/16\le\lambda-3.\label{eq:terminal_soft_update}
    \end{align}
    This is an explicit-budget auxiliary to
    \cite[Figure~\texttt{fig:algo1}, line 1631, and lines 1793--1807]{Arad2016RigorousRG}.
    It concerns the last untrimmed image, not the final collectively
    trimmed space. The physical calibration and the successful complete
    parent event must supply the two comparison bounds.
\end{lemma}
\begin{proof}\leanok
    \uses{lem:native_soft_parent_energy}
    The actual regional construction gives the positive operator
    $C=R-e_{\mathrm s}I$, with norm at most $A_*$ and energy at most
    $2+\rho$ on $G$. The first comparison transfers this energy to $T$.
    Projecting $T$ onto $Y\otimes\mathcal H_{J^c}$ then gives energy
    at most three under the stated budgets. The target has positive
    dimension, and the second loss is less than one, so some projected
    vector is nonzero. The exact regional tensor identity turns it
    into a tensor trial for $R_{\mathrm{nat}}$ of energy at most
    $e_{\mathrm s}+3$. The compression variational bound and
    $R_{\mathrm{nat}}-e_{\mathrm s}I\ge0$ give the first interval in
    (\ref{eq:terminal_soft_update}). The actual minimum comparison and
    $\rho\le1/16$ give the remaining assertions.
\end{proof}

\begin{lemma}[The calibrated terminal energy estimate]
    \label{thm:calibrated_soft_terminal_compressed_energy}
    \lean{TNLean.MPS.AreaLaw.ALVV.exists_calibrated_native_soft_terminal_compressed_energy_estimate}
    \leanok
    Retain the original chain, intervals and untrimmed image of
    Lemma~\ref{thm:soft_terminal_compressed_energy}. Let $g>0$,
    $m,D\ge2$, $t,T_0\ge156+2g$, and set
    \begin{align}
        M_*  & =3t(1+\log m)+4\ell+34, &
        \tau & =(D^{24}M_*^{10})^{-1}.
    \end{align}
    Suppose $0\le\Delta\le(D^{12}M_*^{10})^{-1}$ and
    \begin{align}
        3\left(32e^{-(T_0-156-2g)/8}
        +\frac tm\left(\frac{T_0}t\right)^m\right)
        \le\frac{\tau g}{10^6}.
    \end{align}
    For each admissible center, use a target of dimension $\dim G$,
    with directed comparison loss $3\tau/200$ from $G$, and the same
    actual image with viability loss $\Delta/(91/100)^2$ for that target.
    The actual regional soft matrix, its attained minimum and its
    compression then satisfy all the conclusions of
    Lemma~\ref{thm:soft_terminal_compressed_energy}, including the
    one-sided update of width $49/16$.
    No upper bound on $g$ is imposed. The scalar assumptions are the
    same calibrated scales; the target and image certificates must be
    supplied by that same physical family and complete parent event.
\end{lemma}
\begin{proof}\leanok
    \uses{thm:soft_terminal_compressed_energy}
    The two norm formulas give $0\le A_*\le M_*$, $3t\le M_*$ and
    $M_*\ge34$. In particular $g\le M_*$. The calibrated denominators
    give $\tau M_*\le1/64$,
    $3\tau/200\le1/[1024(A_*+1)]$ and
    $\Delta/(91/100)^2\le1/[100(A_*+1)]$.
    The exponential with shift $48$ is at most the one with shift
    $156+2g$, and $\tau g\le1$, so the middle residual is at most
    $1/16$. Apply the explicit-budget theorem to the same target,
    family image, matrix and regional minima.
\end{proof}

\begin{lemma}[An admissible middle center from the actual children]
    \label{thm:regional_middle_center_from_children}
    \lean{TNLean.MPS.AreaLaw.ALVV.exists_regional_middle_center_of_child_estimates}
    \leanok
    Retain the original interaction assumptions and geometry
    $l\le u\le v\le r\le L$. Suppose the actual attained parent and
    middle minima $e_P,e_M$ are already specified with unit eigenvectors
    and positive scalar shifts. There are actual attained collar minima
    $e_{\mathrm L},e_{\mathrm R}$, with unit eigenvectors and positive
    shifts on $[l,u)$ and $[v,r)$, chosen before the child estimates.

    For every $l\le k\le r$, suppose $p_1,p_2$ are one-sided lower
    estimates of the actual minima on $[l,k)$ and $[k,r)$, each with
    error at most $49/16$. Then
    \begin{align}
        e_P-57/8 & \le p_1+p_2\le e_P,                               \\
        e_M-10   & \le p_1+p_2-e_{\mathrm L}-e_{\mathrm R}-2\le e_M.
    \end{align}
    Both estimates retain attained eigenvector and positive-shift
    certificates for their original parent and middle Hamiltonians.
    Thus the same child-dependent center is admissible in
    Lemma~\ref{thm:calibrated_soft_terminal_compressed_energy}.
    This is a regional auxiliary to
    \cite[Figure~\texttt{fig:algo1}, line 1631, and line 1807]{Arad2016RigorousRG}.
\end{lemma}
\begin{proof}\leanok
    \uses{lem:regional_central_energy_estimate,lem:adjacent_interval_energy_merge}
    The original parent minimum lies between the sum of the two child
    minima and that sum plus one. Adding the two child errors therefore
    gives $1+2(49/16)=57/8$. The actual three-piece regional decomposition
    places $e_P$ between
    $e_{\mathrm L}+e_M+e_{\mathrm R}$ and that sum plus two.
    Apply its middle-center estimate with the parent-sum width
    $57/8\le8$. Uniqueness of attained minima identifies the parent and
    middle values with the already specified $e_P,e_M$.
\end{proof}

\begin{lemma}[A terminal node update from the actual child estimates]
    \label{thm:calibrated_soft_terminal_node_update}
    \lean{TNLean.MPS.AreaLaw.ALVV.exists_calibrated_soft_terminal_node_energy_update}
    \leanok
    Retain the primitive interaction assumptions, arbitrary retained interval
    and calibrated scalar hypotheses of
    Lemma~\ref{thm:calibrated_soft_terminal_compressed_energy}.
    The actual attained parent and middle minima $e_P,e_M$ and the actual
    collar minima $e_{\mathrm L},e_{\mathrm R}$ can all be chosen before
    the child split and its energy estimates, with unit eigenvectors and
    positive scalar shifts for the original regional Hamiltonians.

    For every $l\le k\le r$ and one-sided child estimates $p_1,p_2$ of
    width $49/16$, form
    \begin{align}
        c   & =p_1+p_2-e_{\mathrm L}-e_{\mathrm R}-2, \\
        R_c & =P-M+cI+h_{m,t}(M-cI).
    \end{align}
    Then $c\in[e_M-10,e_M]$, with an attained original-middle energy
    certificate. Define the native soft matrix by the same fixed-complement
    restriction of $R_c$. Its actual soft minimum precedes every target
    and terminal image.

    For each target and the same untrimmed family image satisfying the
    comparison and exact viability certificates of
    Lemma~\ref{thm:calibrated_soft_terminal_compressed_energy}, the actual
    compression has an ordered orthonormal eigenbasis and least eigenvalue
    $\lambda$ satisfying
    \begin{align}
        e_{\mathrm s}\le\lambda\le e_{\mathrm s}+3,\qquad
        e_P-49/16\le\lambda-3\le e_P.
    \end{align}
    Thus $\lambda-3$ is an original-parent lower energy estimate of the
    same width as the children. The original parent eigenvector and positive
    scalar shift have been retained; no new family, terminal matrix or
    candidate is selected. The image certificate must come from the actual
    successful repeated-sampling history, rather than final retained-space
    viability alone. This is an auxiliary node estimate for
    \cite[Figure~\texttt{fig:algo1}, line 1631, and line 1807]{Arad2016RigorousRG}.
\end{lemma}
\begin{proof}\leanok
    \uses{thm:calibrated_soft_terminal_compressed_energy,
        thm:regional_middle_center_from_children}
    Choose the actual regional minima once using the calibrated terminal
    theorem. Apply the child-center theorem to those same parent and middle
    certificates; it chooses the collars before all child estimates and
    proves admissibility of $c$. Apply the retained terminal theorem at
    this derived center to the same target and untrimmed image.
\end{proof}


\begin{lemma}[Energy estimation on a selected interval family]
    \label{lem:selected_soft_terminal_node_energy}
    \lean{TNLean.MPS.AreaLaw.ALVV.exists_selected_soft_terminal_node_energy_update}
    \leanok
    Retain the interaction, geometry and numerical assumptions of
    Lemma~\ref{thm:calibrated_soft_terminal_node_update}. Let $H$ be the
    original chain Hamiltonian, let $\mu\in\operatorname{spec}(H)$, and
    suppose every spectral value is at least $\mu$. Write $G_\mu$ for its
    full ground space. Fix a family $A_i$ on the retained interval and an
    attained middle minimum $e_2$, with a unit vector $x_2$ satisfying
    $Mx_2=e_2x_2$ and $M-e_2I\ge0$.

    The original parent and collar minima can be chosen before every child
    split and its one-sided estimates $p_1,p_2$ of width $49/16$.
    The resulting center $c=p_1+p_2-e_{\mathrm L}-e_{\mathrm R}-2$ lies
    in $[e_2-10,e_2]$. The same soft regional operator
    $R_c=P-M+cI+h_{m,t}(M-cI)$ and its native restriction have the
    minimum and norm bounds of
    Lemma~\ref{thm:calibrated_soft_terminal_node_update}.

    For any native coordinate identification $I_e$ and any matrix $Q$,
    retain the actual untrimmed image
    $Y=\sum_i A_i I_e(\operatorname{range}Q)$.
    Suppose $T$ has the dimension of $G_\mu$, is within directed squared
    error $3\tau/200$ from $G_\mu$, and $Y$ is
    $\Delta/(91/100)^2$-viable for $T$ in the canonical interval/exterior
    tensor coordinates. Then the compression of the same native $R_c$
    to $Y$ has an ordered orthonormal eigenbasis and least eigenvalue
    $\lambda$ satisfying
    \begin{align}
        e_{\mathrm s}\le\lambda\le e_{\mathrm s}+3,\qquad
        e_P-49/16\le\lambda-3\le e_P.
    \end{align}
    The family and terminal image are unchanged; the conclusion concerns
    this untrimmed image, rather than the final trimmed output.
\end{lemma}
\begin{proof}\leanok
    \uses{thm:calibrated_soft_terminal_node_update,
        lem:ground_spectrum_minimum,thm:normalized_eigenvector_energies,
        thm:regional_tensor_viability}
    The attained spectral lower bound gives $\lambda_{\min}(H)=\mu$.
    Uniqueness of the attained minimum of $M$ identifies $e_2$ with the
    middle minimum in the terminal node estimate. The exterior-first
    tensor coordinates are the canonical interval coordinates. Apply that
    estimate to the fixed family and the same matrix $Q$.
\end{proof}
